\clearpage
\begin{summarybox}
\textbf{\textcolor{CSummary}{一句话记住}}\quad
\textbf{中心对称：横坐标取平均，纵坐标也取平均。}
\textbf{绕 $(a,b)$ 转半圈不变}，就是输入的平均数为 $a$，输出的平均数为 $b$。
\[
\boxed{\ f(a+t)+f(a-t)=2b\ }.
\]

\textbf{\textcolor{CSummary}{考试三步速查}}
\begin{enumerate}[leftmargin=*,itemsep=0.2em]
\item \textbf{看范围：} $a+t$ 有定义时，$a-t$ 是否也有定义？
\item \textbf{找中间：}输入之和 $2a$；输出之和 $2b$。
\item \textbf{下结论：}所有成对的点都满足，才能说图像关于 $(a,b)$ 中心对称。
\end{enumerate}

\textbf{\textcolor{CSummary}{两种等价的公式}}
\[
\boxed{\ f(x)+f(2a-x)=2b\ },\qquad
\boxed{\ g(t)=f(a+t)-b\text{ 是奇函数}\ }.
\]
移项写成 $f(p+t)=c-f(q-t)$ 时，先把两个函数值相加，读出中心 $\bigl((p+q)/2,c/2\bigr)$；等式须覆盖全部配对，定义域也须关于横坐标 $(p+q)/2$ 成对。

\textbf{\textcolor{CSummary}{三种证明记忆}}\quad
中点法（最容易理解）／对称点坐标法（最快）／平移成奇函数（与旧知识联系）。

\textbf{\textcolor{CSummary}{四个不易错的提醒}}
\begin{itemize}[leftmargin=*,itemsep=0.2em]
\item 轴对称看\textbf{相等}：$f(a+t)=f(a-t)$；点对称看\textbf{和}：$f(a+t)+f(a-t)=2b$。
\item 若 $a\in D$，则 $f(a)=b$；若 $a\notin D$，中心\textbf{未必在图像上}。
\item 配对求和往往无需解析式：每一对的和都是 $2b$。
\item 两个\textbf{同样高}的不同对称中心能推出周期，不同高度则不能直接套用。
\end{itemize}
\end{summarybox}
